% !TEX program = xelatex
%
% 필요 엔진: XeLaTeX 또는 LuaLaTeX (kotex 사용, pdfLaTeX 로는 조판되지 않음)
%   컴파일:  xelatex -interaction=nonstopmode analysis/interim-report-ko.tex
% 파일 인코딩: UTF-8
%
\documentclass[11pt]{article}

\usepackage{kotex}
\usepackage[a4paper,margin=28mm]{geometry}
\usepackage{booktabs}
\usepackage{array}
\usepackage{caption}
\usepackage{microtype}
\usepackage[hidelinks]{hyperref}

\captionsetup{font=small,labelfont=bf,skip=6pt}
\setlength{\parskip}{0.45em}
\setlength{\parindent}{0pt}
\renewcommand{\arraystretch}{1.15}

\newcolumntype{R}{>{\raggedleft\arraybackslash}p{2.1cm}}
\newcolumntype{L}{>{\raggedright\arraybackslash}p{5.4cm}}

\title{\bfseries 도제 기간 궤적 연구\\[0.3em]\large 중간 연구 보고서}
\author{\normalsize \texttt{trajectories} 연구 저장소}
\date{\normalsize 2026년 8월 22일 기준 \\[0.2em]
      \small 조사 인원 160명 / 포함 134행 / revenue-strict 60행 시점의 스냅샷}

\begin{document}
\maketitle

\begin{abstract}
\noindent
이 문서는 진행 중인 연구의 \textbf{중간 보고서}이며, 최종 결과가 아니다.
현재 조사가 끝난 160명을 기준으로 revenue-strict 부분집합 60행에서
교육 종료 시점부터 첫 성공까지의 중앙값(median)은 \textbf{9.2년},
95\% 신뢰구간(confidence interval, CI)은 $[8.5,\ 13.0]$이다.
사전 등록된 중단 규칙은 아직 \texttt{STOP: False}를 반환하며, 이 수치는 수렴하지 않았다.
지금까지의 수치가 시사하는 바를 한 문장으로 적으면,
\textbf{이 방식으로 연구가 계속될 경우 도제 기간은 대략 9년 부근으로 보인다}는 것이다.
다만 이 값은 \textbf{성공한 사람들 사이에서의} 소요 기간이며 분모가 없다.
성공 확률에 대해서는 아무것도 말하지 않는다.
측정된 편향 중 가장 강한 두 가지가 모두 값을 짧게 미는 방향이므로,
9.2년은 대표값이 아니라 \textbf{하한선}으로 읽어야 한다.
\end{abstract}

\section{이 보고서가 다루는 시점}

이 보고서의 모든 수치는 \textbf{조사 인원 160명, 포함 134행, revenue-strict 60행}
시점의 상태다. 파일럿 10명과 웨이브 1$\sim$6(각 25명)이 종료된 상태이며,
\textbf{웨이브 7과 8(p161\,--\,p210)은 현재 추첨되어 조사가 진행 중이지만
이 보고서의 어떤 숫자에도 포함되어 있지 않다.}
작업 중 \texttt{data/anchors.csv}가 160행을 넘어 늘어나더라도
이 문서는 160행 상태를 보고한다는 원칙으로 작성되었다.

이 문서에 실린 모든 수치는 다음 명령의 실제 출력에서 가져온 것이다.

\begin{quote}\ttfamily\small
PYTHONUTF8=1 python -m src.stats analysis/clocks.csv\\
PYTHONUTF8=1 python -m src.stats analysis/clocks.csv -{}-history analysis/wave\_medians.txt
\end{quote}

\section{무엇을 측정하는가}

측정 대상은 \textbf{형식 교육이 끝난 해부터 첫 성공(first hit)이 일어난 해까지의 경과 연수}다.
저장소에서 이 값은 \texttt{clock\_education}이며, 여섯 개의 날짜 앵커
(\texttt{a1} 출생 $\sim$ \texttt{a6} 대규모 성공) 중
\texttt{a2\_education\_end}와 \texttt{a5\_first\_hit}의 차이로 계산된다.

\subsection{판정 기준은 불변 2026년 달러 기준이다}

상업 버킷에서 \texttt{rev10}은 \textbf{명목 1{,}000만 달러가 아니라
불변(constant) 2026년 달러 기준 1{,}000만 달러}를 뜻한다. 이 구분은 장식이 아니다.
명목 기준을 썼다면 1960년에 창업한 사람은 2020년에 창업한 사람보다
실질 기준으로 약 열한 배 큰 사업을 만들어야 동일한 "첫 성공"으로 인정받게 된다.
즉 명목 기준은 오래된 경력의 도제 기간을 일괄적으로 부풀린다.
1960년의 실제 문턱값은 \$919{,}417이며, \texttt{python -m src.cpi \textless{}연도\textgreater{}}로 조회한다.
첫 파일럿에서 이 규칙을 놓친 결과 한 행이 11년, 다른 한 행이 6년 늦게 기록된 적이 있다.

비상업 버킷에는 등가 기준(basis \texttt{equivalent})이 적용된다.
\texttt{science\_research}는 수상 발표 연도(\texttt{prize}),
\texttt{media\_creators}는 자기 작업의 100만 관객 도달(\texttt{aud1m}),
\texttt{investors\_finance}는 펀드매니저의 경우 1억 달러 초과 첫 펀드 결성(\texttt{fund100}),
애널리스트\,·\,이코노미스트의 경우 공인된 업계 순위 1위 도달(\texttt{rank1})이다.
\texttt{trade\_import\_logistics}는 매출이 우선이되, 매출이 존재하지 않는 것이 정상인 업종이므로
제3자가 발행한 물동량 순위의 1위 등재(\texttt{rank1})를 허용한다.

기업공개(IPO)나 5{,}000만 달러 초과 인수(\texttt{acq50})는 \textbf{대체 기준(fallback)}이며,
더 이른 문턱 통과가 있었다고 알려진 경우에는 사용이 금지된다.
근거가 이른 통과를 시사하지만 연도를 특정할 수 없으면 범위 날짜(\texttt{YYYY-YYYY}, 최대 10년 폭)로 기록하고,
범위조차 잡히지 않으면 그 행은 \texttt{crossing\_undatable} 사유로 제외한다.
연도를 추정하지 않는다는 원칙이 이 연구의 가장 강한 규칙이다.

\section{표본 구성}

\subsection{버킷과 비중}

버킷 비중은 \texttt{frame.md}에 동결되어 있으며 최대잉여법으로 웨이브마다 누적 배분된다.
추첨된 160명의 실제 구성은 아래와 같고, 목표 비중과 사실상 일치한다.

\begin{table}[h]
\centering
\caption{버킷별 구성 (추첨 160명 기준, 포함 행은 제외 처리 후 134행 기준)}
\begin{tabular}{L R R R R}
\toprule
버킷 & 목표 비중 & 추첨 인원 & 실제 비중 & 포함 행 \\
\midrule
\texttt{software\_internet}           & 25\,\% & 40 & 25.0\,\% & 36 \\
\texttt{hardware\_deeptech}           & 15\,\% & 24 & 15.0\,\% & 24 \\
\texttt{consumer\_retail\_industrial} & 13\,\% & 21 & 13.1\,\% & 13 \\
\texttt{science\_research}            & 12\,\% & 19 & 11.9\,\% & 19 \\
\texttt{investors\_finance}           & 10\,\% & 16 & 10.0\,\% & 14 \\
\texttt{media\_creators}              & 10\,\% & 16 & 10.0\,\% & 14 \\
\texttt{healthcare\_biotech}          & 10\,\% & 16 & 10.0\,\% & \phantom{0}9 \\
\texttt{trade\_import\_logistics}     & \phantom{0}5\,\% & \phantom{0}8 & \phantom{0}5.0\,\% & \phantom{0}5 \\
\midrule
합계 & 100\,\% & 160 & 100\,\% & 134 \\
\bottomrule
\end{tabular}
\end{table}

\subsection{지명 출처 목록}

자유 회상으로 들어온 이름은 없다. 모든 명단 항목은 자신이 나온 목록을 인용한다.
160명의 출처 분포는 다음과 같다.

\begin{table}[h]
\centering
\caption{출처 목록별 추첨 인원 (160명)}
\begin{tabular}{L R L R}
\toprule
출처 목록 & 인원 & 출처 목록 & 인원 \\
\midrule
Forbes World's Billionaires (self-made) & 36 & Computer History Museum Fellows & 5 \\
Y Combinator Top Companies              & 18 & Hurun Rich List (중국\,·\,인도)  & 8 \\
Fortune 40 Under 40                     & 15 & 노벨상 (물리\,·\,화학\,·\,의학\,·\,경제) & 12 \\
Endeavor Entrepreneur network           & 15 & Fields Medal                    & 3 \\
Midas List / Midas List Europe          & 15 & Time 100                        & 3 \\
Academy Awards                          & \phantom{0}6 & Turing Award            & 2 \\
Forbes 400 (self-made 6\,--\,10)        & \phantom{0}5 & Breakthrough Prize      & 2 \\
Pulitzer Prize                          & \phantom{0}5 & Sunday Times Rich List (영국) & 2 \\
Grammy Awards                           & \phantom{0}5 & Nikkei / Toyo Keizai (일본) & 2 \\
                                        &              & 매일경제 / 조선일보 (한국) & 1 \\
\bottomrule
\end{tabular}
\end{table}

\subsection{비조종(non-steering) 선정 규칙}

웨이브 6에서 확립되어 현재 사용 중인 선정 규칙은 다음과 같다.
\textbf{목록을 발행 순서대로 열거하고, 그 성공 주체(hit entity)가
\texttt{data/anchors.csv}와 \texttt{data/roster.csv} 어디에도 등장하지 않는 항목을 앞에서부터 취한다.}
Y Combinator Top Companies는 공개 API로 기계 열거가 가능해(91개 항목) 이 규칙에 잘 맞는다.

여기서 중요한 것은 \textbf{유명하다는 이유로 이름을 건너뛰지 않는다}는 점이다.
\texttt{docs/BIASES.md} 24번 항목은 웨이브 6의 명단 작성 에이전트가
"너무 뻔하다"는 이유로 Peter Thiel, Vinod Khosla, Jensen Huang 등을 열거된 목록에서 배제한 사건을 기록한다.
유명한 이름을 회상해서 넣는 것이 편향이듯, 유명한 이름을 피하는 것도 정확히 같은 크기의 편향이다.
두 편향은 상쇄되지 않고 분산만 키운다.

\subsection{쿼터에서 공변량으로의 전환, 그리고 웨이브 1의 불연속}

이 표본의 가장 중요한 구조적 사실은 \textbf{두 개의 서로 다른 표집 설계가 섞여 있다}는 것이다.

파일럿과 웨이브 1은 교차 하한선이 강제되던 체제에서 수집되었다.
비미국 $\geq$ 30\,\%, 1995년 이전 성공 $\geq$ 25\,\%, 여성 $\geq$ 20\,\%가 요구되었고,
이를 위반할 웨이브는 조사 전에 손으로 재조정되었다.
그 결과 해당 구간의 구성은 1995년 이전 50\,\%, 비미국 69\,\%, 여성 42\,\%로,
출처 목록이 그냥 내놓았을 값보다 훨씬 높다.

웨이브 1 이후 연구 저자가 하한선을 제거했다.
현재 \textbf{입장 기준은 금전 기준 하나뿐}이며,
국가(\texttt{country\_primary}), 시대(\texttt{era}), 성별(\texttt{gender})은
\textbf{선정 기준이 아니라 기록되는 공변량}이다.
이 구분은 실질적이다. 변수를 기록하면 그 변수가 답을 움직이는지 발견할 수 있지만,
그 변수로 선정하면 손으로 고정해 버린 것이므로 발견 자체가 불가능해진다.
특히 시대는 도제 기간의 진짜 교란 요인이다.
자본 가용성, 시장 규모, 정보 흐름이 모두 인터넷 시대를 전후해 달라졌기 때문에,
표본이 1995년 이후로 완전히 쏠리면 이 연구는 일반적 질문 대신
"2000\,--\,2020년대에는 얼마나 걸렸는가"에 조용히 답하게 된다.

포함된 134행의 현재 공변량 분포는 아래와 같다.

\begin{table}[h]
\centering
\caption{공변량 분포 (포함 134행)}
\begin{tabular}{L R R}
\toprule
구분 & 행 수 & 비중 \\
\midrule
비미국 경력 & 86 & 64.2\,\% \\
미국 경력   & 48 & 35.8\,\% \\
여성        & 45 & 33.6\,\% \\
남성        & 89 & 66.4\,\% \\
1995년 이전 성공 & 28 & 20.9\,\% \\
1995년 이후 성공 & 106 & 79.1\,\% \\
\bottomrule
\end{tabular}
\end{table}

이 불연속은 기울기가 아니라 계단이다.
초기 행들은 이후 행들에 비해 오래되고, 비미국이며, 여성인 경력을 과대표집한다.
따라서 통합 중앙값은 두 시대의 가중 혼합이며,
그 가중치는 세계에 대한 무엇이 아니라 수집 도중의 설계 변경이 정했다.
올바른 처리는 초기 행을 사후에 제거하는 것이 아니라
$N$이 충분해지면 설계 변경 이후 행만으로 중앙값을 따로 내어 통합값과 대조하는 것이다.
초기 행은 잘못 수집된 데이터가 아니라 다른 규칙으로 정확히 수집된 데이터다.

\section{현재 수치}

\subsection{헤드라인}

\begin{table}[h]
\centering
\caption{revenue-strict 헤드라인 (\texttt{hit\_basis = primary}, \texttt{clock\_education})}
\begin{tabular}{L R}
\toprule
항목 & 값 \\
\midrule
revenue-strict $n$ & 60 \\
중앙값(median)     & 9.2년 \\
95\% 신뢰구간      & $[8.5,\ 13.0]$ \\
신뢰구간 반너비(half-width) & 2.25년 \\
\bottomrule
\end{tabular}
\end{table}

신뢰구간은 부트스트랩으로 산출되며 중앙값을 중심으로 비대칭이다.
하한 8.5년은 중앙값에서 0.7년 떨어져 있는 반면 상한 13.0년은 3.8년 떨어져 있다.
즉 데이터가 허용하는 범위는 짧은 쪽보다 긴 쪽으로 훨씬 넓다.

\subsection{정의 엄격도에 따른 비교}

헤드라인이 어떤 정의 위에 서 있는지 보이기 위해 세 수준을 나란히 출력한다.

\begin{table}[h]
\centering
\caption{정의 엄격도별 중앙값 (\texttt{clock\_education})}
\begin{tabular}{L R R}
\toprule
수준 & 중앙값 (95\% CI) & $n$ \\
\midrule
revenue-strict (매출 \$10M 기준만) & 9.2년 $[8.5,\ 13.0]$ & 60 \\
전 상업 버킷 (IPO\,·\,인수 대체 기준 포함) & 10.0년 $[8.5,\ 13.5]$ & 68 \\
통합 (전 버킷, 혼합 정의) & 14.0년 $[10.0,\ 15.0]$ & 103 \\
\bottomrule
\end{tabular}
\end{table}

헤드라인과 통합값의 차이는 \textbf{4.8년}이다.
이 격차의 대부분은 비상업 버킷, 특히 \texttt{science\_research}(통합 기준 29.0년, $n$=19)가 통합값을 끌어올리기 때문이다.
서로 다른 종류의 사건에는 원리적으로 환율이 없으므로,
헤드라인은 정의가 섞이지 않는 revenue-strict 쪽으로 고정되어 있다.

\subsection{보조 시계}

\begin{table}[h]
\centering
\caption{전 버킷 통합 기준 시계 비교}
\begin{tabular}{L R R}
\toprule
시계 & 중앙값 (95\% CI) & $n$ \\
\midrule
\texttt{clock\_education} (교육 종료 $\to$ 첫 성공) & 14.0년 $[10.0,\ 15.0]$ & 103 \\
\texttt{clock\_age18} (만 18세 $\to$ 첫 성공)       & 20.0년 $[18.0,\ 23.0]$ & 121 \\
\texttt{clock\_venture} (첫 창업 $\to$ 첫 성공)     & \phantom{0}5.0년 $[4.0,\ 6.5]$ & 114 \\
\texttt{age\_at\_first\_hit} (첫 성공 시 연령)      & 38.0년 $[36.0,\ 41.0]$ & 121 \\
\bottomrule
\end{tabular}
\end{table}

행 수가 더 많다는 이유로 헤드라인 시계를 바꾸는 것은 금지된다.
그것은 결과를 얻기 위해 측정을 고르는 일이기 때문이다.

\subsection{웨이브별 중앙값 이력}

\begin{table}[h]
\centering
\caption{웨이브 종료 시점의 revenue-strict 중앙값 (\texttt{analysis/wave\_medians.txt})}
\begin{tabular}{L R R R}
\toprule
기록 시점 & 누적 조사 인원 & revenue-strict $n$ & 중앙값 \\
\midrule
웨이브 1 마감 & \phantom{0}35 & 13 & 10.0년 \\
웨이브 2 마감 & \phantom{0}60 & 21 & \phantom{0}9.0년 \\
웨이브 3 마감 & \phantom{0}85 & 28 & \phantom{0}8.8년 \\
중간 기록     & --- & 34 & \phantom{0}9.0년 \\
웨이브 4 마감 & 110 & 45 & \phantom{0}9.0년 \\
웨이브 5 마감 & 135 & 52 & \phantom{0}9.8년 \\
웨이브 6 마감 & 160 & 60 & \phantom{0}9.2년 \\
\bottomrule
\end{tabular}
\end{table}

파일럿 10명 시점은 로그에 남아 있지 않다. 로그의 첫 항목은 웨이브 1 마감 시점이며,
네 번째 항목($n$=34)은 웨이브 마감이 아니라 중간에 기록된 값으로,
중앙값이 $n$=30 하한선을 처음 넘어 출력된 시점이다.
누적 조사 인원 열은 웨이브 구조(파일럿 10명 + 웨이브당 25명)에서 재구성한 것이고,
$n$과 중앙값 두 열은 로그 파일에 있는 그대로다.

중앙값은 기록이 시작된 이래 8.8년에서 10.0년 사이의 좁은 띠 안에서 움직였고 추세를 보이지 않는다.
다만 마지막 두 기록의 이동폭이 각각 0.80년과 0.60년으로,
중단 규칙이 요구하는 0.50년보다 크다.

\section{중단 규칙}

중단 규칙은 웨이브 1 이전에 사전 등록되었고 이후 변경되지 않았다.
\texttt{hit\_basis = primary}인 revenue-strict 부분집합에서
\texttt{clock\_education}의 중앙값과 95\,\% 부트스트랩 신뢰구간을 매 웨이브 계산하고,
다음 두 조건이 \textbf{동시에} 만족될 때 수집을 멈춘다.

\begin{enumerate}
\item 연속 두 웨이브에 걸쳐 중앙값의 이동이 각각 0.5년 미만일 것
\item 부트스트랩 신뢰구간의 반너비가 1.0년 이하일 것
\end{enumerate}

또한 revenue-strict $n$이 30에 이르기 전에는 중앙값을 아예 들여다보지 않는다.
이 하한선은 선택적 중단(optional stopping)을 막기 위한 것이다.
매 웨이브 확인하면서 수치가 안정되어 보이는 순간에 멈추면
잡음이 우연히 작았던 웨이브에서 우선적으로 멈추게 되고,
그 결과 실제보다 정밀해 보이는 중앙값이 나온다.
$n$=30 하한선은 이미 통과했으므로 현재 중앙값은 출력된다.

현재 규칙의 판정은 다음과 같다.

\begin{quote}\ttfamily\small
STOP: False - median drift too large (0.80, 0.60 yr; need both \textless{} 0.50)
\end{quote}

즉 조건 1이 깨져 있다. 조건 2는 더 크게 깨져 있다.
반너비 2.25년은 문턱값 1.00년의 두 배가 넘는다.

\subsection{규칙이 함의하는 표본 크기}

중앙값의 정밀도는 대략 $\sqrt{n}$에 비례해 개선된다.
반너비를 2.25년에서 1.00년으로 줄이려면 $n$이 약 $(2.25/1.00)^2 \approx 5.06$배 필요하므로,
\textbf{revenue-strict 약 300행}이 요구된다.
현재 조사 인원 대비 revenue-strict 산출률이 $60/160 = 37.5\,\%$이므로,
이 비율이 유지된다면 \textbf{조사 인원 약 800명}에 해당한다.
현재 160명이므로 앞으로 여섯 배 이상이 남아 있다는 뜻이다.

이 계산과 관련해 한 가지 불일치를 기록해 둔다.
\texttt{docs/HANDOFF.md}는 같은 취지의 추정을
"revenue-strict $n \approx 230$, 조사 인원 약 500명"으로 적고 있으나,
현재 반너비 2.25년(또는 같은 문서가 다른 절에서 언급하는 2.50년)에서 출발한
위 산술은 그 값을 재현하지 않는다.
이 보고서는 실제 CLI 출력값에서 직접 계산한 300행 / 800명 쪽을 사용한다.

수렴하지 않았다는 사실 자체는 규칙이 정상 작동하고 있다는 신호다.
다만 사전 등록된 규칙이 아직 만족되지 않은 상태에서
현재 중앙값과 그 실제 구간, 그리고 수렴하지 않았다는 진술을 함께 공표하는 것 역시
정당한 과학적 산출물이다. 어느 쪽을 택할지는 연구 저자의 결정 사항이다.

\section{데이터 품질}

\subsection{제외 감사}

\textbf{160행 중 26행(16.2\,\%)이 제외되었다.}
제외된 행은 다른 어떤 표에도 나타나지 않으므로,
무엇이 제외되는지에 편향이 있다면 오직 이 절에서만 보인다.

\begin{table}[h]
\centering
\caption{제외 사유}
\begin{tabular}{L R}
\toprule
사유 & 행 수 \\
\midrule
\texttt{crossing\_undatable} (문턱 통과 연도를 특정도 범위 지정도 못 함) & 22 \\
\texttt{roster\_unverified} (명단 출처가 그 인물을 실제로 담고 있지 않음) & \phantom{0}3 \\
\texttt{criterion\_never\_met} (기준을 끝내 충족하지 못함) & \phantom{0}1 \\
\midrule
합계 & 26 \\
\bottomrule
\end{tabular}
\end{table}

\begin{table}[h]
\centering
\caption{교차 구분별 제외율}
\begin{tabular}{L R R R}
\toprule
구분 & 포함 & 제외 & 제외율 \\
\midrule
미국   & 48 & \phantom{0}8 & 14.3\,\% \\
비미국 & 86 & 18 & 17.3\,\% \\
여성   & 45 & \phantom{0}6 & 11.8\,\% \\
남성   & 89 & 20 & 18.3\,\% \\
\midrule
\texttt{hardware\_deeptech}           & 24 & \phantom{0}0 & \phantom{0}0.0\,\% \\
\texttt{science\_research}            & 19 & \phantom{0}0 & \phantom{0}0.0\,\% \\
\texttt{software\_internet}           & 36 & \phantom{0}4 & 10.0\,\% \\
\texttt{investors\_finance}           & 14 & \phantom{0}2 & 12.5\,\% \\
\texttt{media\_creators}              & 14 & \phantom{0}2 & 12.5\,\% \\
\texttt{trade\_import\_logistics}     & \phantom{0}5 & \phantom{0}3 & 37.5\,\% \\
\texttt{consumer\_retail\_industrial} & 13 & \phantom{0}8 & 38.1\,\% \\
\texttt{healthcare\_biotech}          & \phantom{0}9 & \phantom{0}7 & 43.8\,\% \\
\bottomrule
\end{tabular}
\end{table}

제외는 흩어진 개별 행의 문제가 아니라 \textbf{버킷 단위로 몰린다.}
\texttt{hardware\_deeptech}와 \texttt{science\_research}는 제외율이 0\,\%인 반면
\texttt{healthcare\_biotech}는 43.8\,\%, \texttt{consumer\_retail\_industrial}은 38.1\,\%다.
이는 무작위 결측이 아니라 해당 업종의 공시 관행이 만드는 구조적 결측이다.
시대(era)는 성공 연도에서 파생되는데 제외된 행에는 대개 그 연도가 없으므로 이 표에서 감사하지 않는다.
제외된 행의 시대를 판정하는 것은 추론이지 측정이 아니다.

\subsection{범위 날짜의 비중}

포함된 134행 중 \textbf{64행(47.8\,\%)이 범위로 지정된 \texttt{a5}}를 가진다.
즉 성공 연도가 \texttt{YYYY-YYYY} 형태로 기록되어 있고, 시계는 그 중점을 사용한다.
범위 날짜는 추측이 아니라 실제 측정이다.
근거가 사건을 두 해 사이에 가두었다는 사실 자체가 정보이며,
마음에 드는 연도를 포함할 때까지 범위를 넓히는 것은 허용되지 않는다(최대 폭 10년).

절반에 가까운 행이 범위 날짜를 갖기 때문에 범위 민감도 분석이 중요하다.
전 버킷 통합 기준의 결과는 다음과 같다.

\begin{table}[h]
\centering
\caption{범위 날짜 민감도 (전 버킷 통합, \texttt{clock\_education})}
\begin{tabular}{L R R}
\toprule
처리 방식 & 중앙값 (95\% CI) & $n$ \\
\midrule
범위 행 제외 & 16.0년 $[13.0,\ 19.0]$ & \phantom{0}54 \\
전체, 범위 중점 사용 & 14.0년 $[10.0,\ 15.0]$ & 103 \\
전체, 가능한 최이른 해 & 12.0년 $[9.0,\ 14.0]$ & 103 \\
전체, 가능한 최늦은 해 & 14.0년 $[12.0,\ 18.0]$ & 103 \\
\bottomrule
\end{tabular}
\end{table}

이 표는 통합 중앙값을 한정할 뿐 revenue-strict 헤드라인을 한정하지 않는다는 점에 주의한다.
같은 이유로 신뢰도 민감도 분석 역시 통합 기준이다.
고신뢰 행만 쓰면 21.0년 $[14.0,\ 27.5]$, $n$=38이고 전체 행은 14.0년 $[10.0,\ 15.0]$, $n$=103으로
\textbf{7.0년 벌어진다}. 포함 134행 중 96행이 고신뢰에 못 미친다.
이 격차는 통합 중앙값을 점추정치로 읽으면 안 된다는 뜻이다.

\subsection{현재의 구속 조건은 \texttt{a2}다}

\textbf{포함된 134행 중 31행의 \texttt{a2\_education\_end}가 \texttt{unknown}이다.}
헤드라인 시계는 \texttt{a2}에서 출발하므로, 이 31행은 헤드라인에 전혀 기여하지 못한다.

구체적으로 보면 이 제약의 크기가 드러난다.
포함 행 중 \texttt{hit\_basis = primary}인 행은 77행인데,
그중 \texttt{a2}가 채워진 행은 60행뿐이다.
즉 \textbf{매출 기준으로 성공 연도가 확정된 17개의 행이 오직 교육 종료 연도가 없다는 이유로
헤드라인 표본에서 빠져 있다.}
이는 \texttt{a5}보다 큰 손실이며, 따라서 \textbf{현재 이 연구의 구속 조건은 \texttt{a5}가 아니라 \texttt{a2}다.}

그리고 이 결측은 무작위가 아니다.
"이 사람이 학교를 언제 마쳤는가"는 미국\,·\,영국 인물에 대해서는 잘 기록되어 있고,
1970\,--\,80년대 인도\,·\,중국 창업자나 보고할 만한 학위를 거치지 않은 경력에 대해서는 잘 기록되어 있지 않다.
선정 기준은 금전 하나뿐인데도, 교육 시계는 조용히
"학력이 기록된 사람들", 즉 서구\,·\,고학력\,·\,최근 경력 쪽으로 조건화된다.

\section{한계}

\texttt{docs/BIASES.md}는 24개의 편향을 방향과 함께 기록한다.
중앙값을 인용하기 전에 반드시 읽어야 할 항목들을 방향과 함께 정리하면 다음과 같다.

\begin{table}[h]
\centering
\caption{주요 편향과 그 방향}
\begin{tabular}{p{0.8cm} p{7.6cm} p{5.6cm}}
\toprule
번호 & 내용 & 중앙값에 미치는 방향 \\
\midrule
1 & 생존 편향. 표본 전체가 결과로 조건화되어 있다. 12년을 일하고 끝내 넘지 못한 사람은 어떤 목록도 표집하지 않는다. & 값의 편향이 아니라 \textbf{의미의 한계}. 데이터를 늘려도 교정 불가. \\
2 & 극단적 성과로 선정하고 완만한 이정표로 날짜를 매긴다. Bezos급에 도달한 사람은 \$10M에서 멈춘 사람보다 그 문턱을 더 빨리 넘었을 가능성이 크고, 후자는 표본에 들어오지 못한다. & \textbf{짧게} \\
3 & \texttt{primary}로의 진입이 무작위가 아니다. 매출을 일찍 공개한 기업, 즉 빨리 성장한 기업이 과대표집된다. & \textbf{짧게} \\
5 & 대체 기준(IPO\,·\,인수)의 지연. 규칙 개정으로 대부분 막혔으나 잔여 지연이 남는다. 포함 행 중 8행이 대체 기준이다. & 해당 행에 대해 길게 \\
17 & 교육 시계가 학력이 기록된 사람에게 조건화된다(위 6.3절). 학력 기록이 없는 창업자는 더 일찍 벌기 시작했을 개연성이 있다. & 부호 불명, 아마 \textbf{짧게} \\
18 & 단일 패스 조사의 측정된 누락률. 제외된 14행을 재조사했더니 \texttt{crossing\_undatable} 13행 중 3행이 실제로는 날짜 지정 가능했다. 즉 \textbf{제외 4건 중 1건꼴로 회수 가능}했고, 회수된 행은 비미국\,·\,1995년 이전으로 쏠린다. & 구성 편향. 남아 있는 제외는 보이는 것보다 약한 근거다. \\
22 & 블라인드 감사의 격리가 환경이 아니라 지시로 유지된다. 웨이브 6에서 4개의 병렬 조사 에이전트가 완성된 답이 담긴 54개 산출물을 공유 작업 디렉터리에 남겼다. 손으로 옮긴 뒤에야 감사가 진행되었다. & 실패할 경우 \textbf{신뢰도를 과대평가}시킨다. \\
23 & 같은 회사 공동창업자는 독립 관측이 아니다. p70 Dubinsky와 p94 Hawkins는 둘 다 1995년 Palm Computing에 걸린다. 한 사건, 두 행이다. & 중앙값에는 영향이 작으나 \textbf{신뢰구간을 부당하게 좁힌다}. 중단 규칙을 직접 공격하는 유일한 실패 모드다. \\
24 & 유명세를 피하는 것도 조종이다. 열거된 목록에서 "너무 뻔하다"는 이유로 이름을 배제한 사례가 기록되어 있다. & 방향 불명. 20번과 반대 방향이며 상쇄되지 않고 분산만 키운다. \\
\bottomrule
\end{tabular}
\end{table}

\subsection{분모가 없다}

이 한계는 다른 모든 것 위에 있으므로 따로 적는다.

중앙값 9.2년은 \textbf{성공한 사람들 사이에서의} 소요 기간이다.
버킷은 Forbes, Midas, Nobel, Hurun처럼 극단적 결과로 선정하는 목록에서 채워지므로,
12년을 일하고 끝내 문턱을 넘지 못한 사람은 이 표본에 한 명도 없다.
그런 사람들은 막대한 수로 존재하며, 드물어서가 아니라 \textbf{구성상} 부재한다.

따라서 이 수치가 뒷받침하는 문장은 다음 하나뿐이다.

\begin{quote}
결국 성공한 사람들 중 가운데 사람은 약 9년이 걸렸다.
\end{quote}

이 수치가 뒷받침하지 \textbf{않는} 문장은 다음과 같다.

\begin{quote}
지금 시작하면 약 9년 안에 성공할 확률이 대략 50\,\%다.
\end{quote}

두 번째 문장에는 분모, 즉 시작한 사람 전부가 필요하고 이 데이터셋에는 분모가 전혀 없다.
복권 당첨자들이 표를 얼마나 오래 들고 있었는지 재는 것과 같다.
정확히 계산된 진짜 숫자이지만 복권을 사는 것이 좋은 계획인지에 대해서는 아무것도 말하지 않는다.

그리고 측정된 편향 중 가장 강한 두 가지(2번과 3번)가 모두 값을 짧게 미는 방향이다.
따라서 정상적으로 성공한 경력의 참 중앙값은
이 연구가 보고하는 값보다 \textbf{길다}고 보는 것이 작업 가정이며,
9.2년은 대표값이 아니라 \textbf{하한선}이다.

\subsection{이 문서가 하지 않는 것}

설계 명세(\texttt{docs/superpowers/specs/2026-07-31-apprenticeship-trajectories-design.md})는
\textbf{완성된 숫자를 저자 자신의 일정에 적용하는 것}을 명시적 비목표(non-goal)로 규정한다.
그것은 숫자가 존재한 뒤에야 시작할 수 있는 별개의 프로젝트다.
성공 확률, 기저율, 실패 건수 추정 역시 비목표다.
따라서 이 보고서에는 제언, 시사점, 조언에 해당하는 절이 없다.
이 문서는 무엇을 어떻게 측정했고, 얼마나 불확실하며, 무엇이 그 값을 바꿀 수 있는지만 기술한다.

\section{향후 계획}

\begin{enumerate}
\item \textbf{웨이브 7\,·\,8 종료.} p161\,--\,p210이 추첨되어 조사 중이다.
      종료 시 조사 인원 210명이 되며, 현재 산출률이 유지되면 revenue-strict는 약 79행이 된다.
      두 웨이브 모두 검증 0오류와 블라인드 감사를 통과해야 마감된다.

\item \textbf{\texttt{a2} 표적 재조사.} 포함 행 중 31행의 교육 종료 연도가 없고,
      그중 17행은 이미 매출 기준 성공 연도가 확정된 \texttt{primary} 행이다.
      이 17행을 채우면 새 인물을 한 명도 추가하지 않고 헤드라인 표본이 60행에서 최대 77행으로 늘어난다.
      단위 노력 대비 표본 증가가 가장 큰 작업이다.

\item \textbf{제외 행 2차 패스.} 편향 18번이 제외 4건 중 1건꼴의 회수 가능성을 측정했고,
      회수는 비미국\,·\,비영어권 근거로 쏠린다.
      최종 중앙값을 내기 전에 22개의 \texttt{crossing\_undatable} 행에 대한 독립 2차 패스가 필요하다.
      단, 회수 패스도 1차 패스와 동일한 회의주의를 적용해야 한다.
      실제로 한 회수 후보(Raghuram Rajan에게 제시된 \textit{Economist} 게스트 네트워크 설문)는
      공인된 업계 순위가 아니므로 올바르게 기각되었다.

\item \textbf{Palm 쌍에 대한 결정.} p70과 p94는 하나의 사건에 걸린 두 행이다.
      상관된 행은 정보를 더하지 않으면서 부트스트랩 구간을 좁히므로,
      연구가 스스로 벌지 않은 정밀도로 자기 중단 조건을 만족시킬 수 있다.
      앞으로의 프레임은 인물당 1행이 아니라 \textbf{성공 사건(주체 + 연도)당 1행}을 요구해야 한다.

\item \textbf{블라인드 감사의 환경적 격리.} 병렬 1차 조사 에이전트는 문제를 키운다.
      에이전트마다 출력 파일을 분리하고, 감사자를 파견하기 전에 1차 산출물을 공유 디렉터리 밖으로 옮겨야 한다.
      또한 \texttt{docs/RESEARCH-RULES.md}가 모든 에이전트에게 \texttt{data/anchors.csv}를 열라고 지시하므로,
      감사 프롬프트에서 이 지시를 명시적으로 배제해야 한다.

\item \textbf{두 버킷의 표집 목록 확보.} \texttt{investors\_finance}의 애널리스트\,·\,이코노미스트 절반과
      \texttt{trade\_import\_logistics} 전체에는 전용 표집 목록이 없다(편향 12번).
      다만 프레임 변경은 그 자체가 편향이므로 웨이브 도중에는 하지 않는다.

\item \textbf{설계 변경 이후 행의 분리 보고.} $N$이 충분해지면
      웨이브 2 이후 행만으로 중앙값을 산출해 통합값과 대조한다(편향 14번).
\end{enumerate}

\vspace{1em}
\hrule
\vspace{0.6em}
{\small
이 보고서의 모든 수치는 커밋 95a1a41(웨이브 6 마감) 시점의
\texttt{data/anchors.csv} 160행, \texttt{analysis/clocks.csv},
\texttt{analysis/wave\_medians.txt}에서 직접 산출되었다.
어떤 수치도 다른 문서에서 옮겨 적지 않았다.
}

\end{document}
